\section{Implications of Execution-Edit Safety}
% 执行编辑安全的含义
\label{sec:implications}

This section explains six consequences of our formal results for runtime design.
% 本节说明我们的形式化结果对运行时设计产生的六个结论。
Here, a \emph{result} means one way the task can finish, such as completing the laptop purchase through supplier \(L\) or \(R\).
% 这里的\emph{结果}是任务的一种完成方式，例如通过供应商 \(L\) 或 \(R\) 完成笔记本电脑采购。
It is not a value returned by a tool.
% 它不是工具返回值。
A result remains required until the execution record shows that it has finished or that the party requiring it allowed its removal.
% 一个结果会持续受到要求，直到执行记录表明它已经完成，或提出该要求的一方允许移除它。

\paragraph{A smaller workflow can be less safe}
% 更小的工作流反而可能更不安全。
In the laptop example, the current workflow allows the school to buy through supplier \(L\) or supplier \(R\).
% 在笔记本电脑例子中，当前工作流允许学校通过供应商 \(L\) 或供应商 \(R\) 完成采购。
Suppose the Agent edits it to keep only \(L\).
% 假设 Agent 编辑该工作流，只保留 \(L\)。
Every remaining call may obey the purchase policy, but the edit removes the still-required \(R\) result.
% 剩余的每个调用都可能遵守采购策略，但该编辑删除了仍然受到要求的 \(R\) 结果。
The edit is safe only if the record shows that the \(R\) purchase has finished or that the school allowed its removal.
% 只有当记录表明通过 \(R\) 的采购已经完成，或学校允许移除它时，该编辑才是安全的。
The runtime must therefore derive what the edit must preserve from the execution record, not from the proposed workflow.
% 因此，运行时必须根据执行记录而不是提议的工作流推导该编辑必须保留什么。

\paragraph{Two safe plans may have no safe way to run together}
% 两个分别安全的计划可能无法安全地共同运行。
Suppose a release workflow must support both normal release and emergency release.
% 假设一个发布工作流必须同时支持普通发布和紧急发布。
The normal plan obtains approval \(y\) before deployment \(x\), producing \(yx\).
% 普通计划先取得审批 \(y\)，再执行发布 \(x\)，形成 \(yx\)。
The emergency plan deploys first and then files an emergency report \(z\), producing \(xz\).
% 紧急计划先执行发布，再提交紧急报告 \(z\)，形成 \(xz\)。
The policy \(\mathcal A=\{\epsilon,x,y,xz,yx\}\) permits both orders, so either plan is safe when checked alone.
% 策略 \(\mathcal A=\{\epsilon,x,y,xz,yx\}\) 允许这两种次序，因此单独检查时，每个计划都是安全的。
In the combined workflow \(X\otimes(Y\oplus Z)\), both plans contain \(x\).
% 在组合工作流 \(X\otimes(Y\oplus Z)\) 中，两个计划都包含 \(x\)。
Allowing \(x\) first leaves the normal plan able to finish only in the forbidden order \(xy\), while blocking \(x\) makes the emergency plan impossible.
% 先允许 \(x\) 会使普通计划只能按被禁止的次序 \(xy\) 完成，而阻止 \(x\) 又会使紧急计划无法执行。
The runtime cannot make either choice while keeping both plans safe, so the exact checker returns \(\mathsf{Reject}\).
% 运行时无论如何选择都无法保证两个计划安全，因此精确检查器返回 \(\mathsf{Reject}\)。

\paragraph{An edit can change from Accepted to Reject without a new authorization}
% 即使没有新授权，编辑也可能从 Accepted 变为 Reject。
An edit verdict applies only to the execution record that the checker examined.
% 编辑判定只适用于检查器所检查的那份执行记录。
In \cref{fig:two-histories}, Restore copies an approval call as \(x_a'\).
% 在 \cref{fig:two-histories} 中，Restore 把一个审批调用复制为 \(x_a'\)。
The original call and \(x_a'\) refer to the same authorized action \(d_a\).
% 原调用和 \(x_a'\) 指向同一个已授权动作 \(d_a\)。
Processing \(x_a'\) therefore adds no authorization-log entry, but it still moves the restored branch past approval.
% 因此，处理 \(x_a'\) 不会增加授权日志记录，但仍会使恢复后的分支越过审批。
A concurrent Merge or Restore that was Accepted before this call may be rejected afterward.
% 在该调用之前被接受的并发 Merge 或 Restore，可能在该调用之后被拒绝。
If the new verdict is Reject, the edit must not take effect in that execution state.
% 如果新判定是 Reject，该编辑就不能在该执行状态下生效。
Requesting authorization for \(d_a\) again cannot undo this progress because the runtime reuses \(d_a\)'s existing authorization record.
% 再次请求对 \(d_a\) 的授权无法撤销这一进展，因为运行时会复用 \(d_a\) 已有的授权记录。
The runtime must discard the earlier verdict and check the edit against the new record.
% 运行时必须丢弃先前判定，并根据新记录重新检查编辑。
A later check can change its answer after a relevant fact changes, such as completing a required result or authorizing its removal, but not by resetting recorded history.
% 在相关事实发生变化后，例如某个要求结果已经完成或其移除获得授权，后续检查可以改变答案，但不能通过重置已记录历史来改变答案。

\paragraph{Saving every call and authorization is still not enough}
% 保存每个调用和授权仍然不够。
Suppose the runtime retains all workflow calls, actions, and authorization records but loses the links among them.
% 假设运行时保留所有工作流调用、动作和授权记录，却丢失它们之间的对应关系。
It can no longer tell whether a restored ``approve purchase 42'' call reuses purchase 42's authorization or needs a new one.
% 它将无法判断恢复出的“批准第 42 号采购”调用是复用第 42 号采购的授权，还是需要新授权。
The tool name and call text cannot recover this information.
% 工具名称和调用文本无法恢复这一信息。
\Cref{thm:necessary-state} proves that two records containing the same calls, actions, and authorizations can require opposite verdicts when these links differ.
% \Cref{thm:necessary-state} 证明，两份包含相同调用、动作和授权的记录，在这些对应关系不同时可能要求相反判定。
An exact runtime must therefore record which action each workflow call and authorization refers to.
% 因此，精确运行时必须记录每个工作流调用和每项授权指向哪个动作。

\paragraph{There is one largest safe set of call orders}
% 存在一个最大的安全调用次序集合。
Safety could in principle admit several incomparable rule sets, each allowing call orders that the others block.
% 原则上，安全性可能允许多个不可比较的规则集合，每个集合都允许其他集合所阻止的调用次序。
\Cref{sec:model} instead proves that the checker returns one safe execution set containing every other safe execution set.
% \Cref{sec:model} 则证明，检查器返回一个包含所有其他安全执行集合的安全执行集合。
The generated runtime rules therefore combine all call orders that can coexist without violating policy or preventing a required result from finishing.
% 因此，生成的运行时规则会合并所有能够共存且不会违反策略、也不会阻止要求结果完成的调用次序。

\paragraph{Exactness survives later calls, edits, and restarts}
% Exactness 会在后续调用、编辑和重启之后继续成立。
The exact guarantee is not limited to the first accepted edit.
% 精确保证并不局限于第一次被接受的编辑。
After allowed calls have run, checking the remaining workflow again gives the same safe ways to finish as the installed rules.
% 在获准调用执行后，再次检查剩余工作流会得到与已安装规则相同的安全完成方式。
Later Fork, Restore, and Merge requests, rule updates, stops, and restarts preserve the same guarantee.
% 后续 Fork、Restore、Merge 请求、规则更新、停止及重启都会保持同一保证。
Every later edit is therefore checked under the same exact safety condition as the first.
% 因此，每个后续编辑都会按照与第一次相同的精确安全条件进行检查。

\Cref{sec:model} formalizes the execution record and safety goal.
% \Cref{sec:model} 形式化定义执行记录和安全目标。
\Cref{sec:exact-check} defines the checker, and \cref{sec:enforcement,sec:results} show how the runtime installs its rules atomically and preserves the guarantees above.
% \Cref{sec:exact-check} 定义检查器，\cref{sec:enforcement,sec:results} 说明运行时如何原子地安装规则并保持上述保证。
